\chapter{Приложение к вариационным неравенствам}\label{ch:viqa}

\section{Постановка и место в работе}\label{sec:viqa_setup}

Главы~\ref{ch:sp-broyden-theory} и~\ref{ch:symmetric} развивают
\emph{локальную} теорию: они количественно ограничивают ошибку
аппроксимации якобиана $\Norm{B_k - J(x_k)}$, но скорость самого
итерационного процесса гарантируется лишь в~окрестности решения. В
настоящей главе показано, куда эта ошибка входит в~\emph{глобальном}
методе второго порядка. Изложение следует
работе~\cite{viqa2024} и сжато до~интерфейса с~главами~\ref{ch:sp-broyden-theory}
и~\ref{ch:symmetric}.

\paragraph{Происхождение результатов.} Метод VIJI, оценка
скорости~\eqref{eq:viqa_rate}, нижняя оценка~\eqref{eq:viqa_lower},
рестарт-схема~\eqref{eq:viqa_restart_rate} и~априорные оценки уровня
неточности~$\delta$ для~L-Broyden получены автором настоящей работы
в~соавторстве~\cite{viqa2024} и~излагаются здесь как внешний метод
второго порядка. Собственный вклад настоящей главы~--- подстановка
мультисекущей оценки глав~\ref{ch:sp-broyden-theory}--\ref{ch:symmetric}
в~роль~$\delta$ (утверждение~\ref{prop:viqa_subst}) и~связанное с~ней
разграничение подпространственного ($O(\rho_k)$) и~полного ($\Fro{E_k}$)
контроля ошибки аппроксимации якобиана.


Рассматривается вариационное неравенство (ВН) Минти~\cite{facchinei2003}: найти
$x^*\in\mathcal{X}$ такое, что
\begin{equation}\label{eq:mvi}
  \langle F(x),\, x - x^*\rangle \ge 0 \qquad \forall\, x\in\mathcal{X},
\end{equation}
где оператор $F\colon\mathcal{X}\to\RR^n$ предполагается
\emph{монотонным}, $\langle F(x)-F(y),\,x-y\rangle\ge 0$, и~$L_1$-гладким,
$\Norm{\nabla F(x)-\nabla F(y)}_{\mathrm{op}}\le L_1\Norm{x-y}$.
Класс~\eqref{eq:mvi} охватывает как минимизацию ($F=\nabla f$), так и
минимаксные задачи $\min_x\max_y\Phi(x,y)$ с~оператором
$F(z)=(\nabla_x\Phi,\,-\nabla_y\Phi)$. Методы второго порядка для~\eqref{eq:mvi}
сходятся быстрее методов первого порядка, но~требуют на~каждом шаге
якобиан $\nabla F$, вычисление которого дорого; замена точного якобиана
аппроксимацией $J(v)\approx\nabla F(v)$ анализируется ниже.

Квазиньютоновские аппроксимации в~методах для~вариационных неравенств
и~монотонных включений~--- активно развиваемое направление: Jiang
и~Mokhtari~\cite{jiang2023accel} построили ускоренный квазиньютоновский
проксимальный экстраградиент с~глобальной неасимптотической
сверхлинейной скоростью, а~Jiang, Jin и~Mokhtari~\cite{jiang2023online}~---
квазиньютоновский метод с~глобальной сверхлинейной сходимостью на~основе
онлайн-обучения кривизны. Настоящая глава рассматривает дополнительный
аспект~--- влияние \emph{неточности} якобиана на~скорость метода
для~ВН~\cite{viqa2024}.

\section{Дуальная экстраполяция второго порядка}\label{sec:viqa_method}

Метод VIJI~\cite{viqa2024} строится по~схеме дуальной экстраполяции и
на~каждой итерации решает вспомогательное ВН с~линейной по~$x$ моделью
оператора
\begin{equation}\label{eq:viqa_model}
  \Psi_v(x) = F(v) + J(v)\,[x-v],
\end{equation}
в~которой $J(v)$~--- приближение якобиана $\nabla F(v)$, играющее ту же
роль, что и~матрица $B_k$ в~главах~\ref{ch:sp-broyden-theory}
и~\ref{ch:symmetric}. Неточность модели~\eqref{eq:viqa_model}
контролируется одним числом:
\begin{equation}\label{eq:viqa_delta}
  \Norm{\nabla F(v) - J(v)}_{\mathrm{op}} \le \delta .
\end{equation}

Одну итерацию метода~\cite{viqa2024} (Алгоритм~1 указанной работы) удобно
записать как четыре шага.
\begin{enumerate}[leftmargin=*]
  \item По~текущему $s_k$ вычисляется
        \[
          v_{k+1}
          = \arg\max_{v\in\mathcal{X}}
            \Bigl\{\langle s_k,\,v-x_0\rangle - \tfrac12\Norm{v-x_0}^2\Bigr\}.
        \]
  \item \emph{Решение подзадачи с~демпфированной моделью.} Находится
        $x_{k+1}\in\mathcal{X}$, приближённо решающее ВН с~оператором
        \begin{equation}\label{eq:viqa_Omega}
          \Omega^{\eta}_{v}(x)
          = \Psi_v(x) + \eta\,\delta\,(x-v) + 5L_1\Norm{x-v}\,(x-v),
        \end{equation}
        где к~линейной модели~\eqref{eq:viqa_model} добавлены два
        регуляризатора: $5L_1\Norm{x-v}(x-v)$
        (стандартное демпфирование) и 
        $\eta\,\delta\,(x-v)$, страхующий от~неточности якобиана. Критерий
        приближённости подзадачи~---
        \begin{equation}\label{eq:viqa_subcrit}
          \sup_{x\in\mathcal{X}}
            \langle\Omega_{v_{k+1}}(x_{k+1}),\,x_{k+1}-x\rangle
          \le \tfrac{L_1}{2}\Norm{x_{k+1}-v_{k+1}}^3
              + \delta\,\Norm{x_{k+1}-v_{k+1}}^2 .
        \end{equation}
  \item \emph{Двойственный шаг.} Подбирается множитель $\lambda_{k+1}>0$
        из~условия
        $\tfrac{1}{32}\le\lambda_{k+1}\bigl(\tfrac{L_1}{2}\Norm{x_{k+1}-v_{k+1}}
        +\beta_{k+1}\bigr)\le\tfrac{1}{22}$.
  \item \emph{Накопление.} $s_{k+1} = s_k - \lambda_{k+1}\,F(x_{k+1})$.
\end{enumerate}
Существенно, что~неточный якобиан входит \emph{только} через
модель~\eqref{eq:viqa_model} внутри~\eqref{eq:viqa_Omega}: вся остальная
механика метода от~$J(v)$ не~зависит, поэтому цена аппроксимации
выражается единственным параметром~$\delta$
из~\eqref{eq:viqa_delta}.

\paragraph{Оценка скорости (Теорема~3.2 из~\cite{viqa2024}).}
\label{par:viqa_rate}
Центральный результат~--- явная зависимость скорости от~уровня
неточности~$\delta$: при $\beta_k=\delta$, $\eta=10$ для
$\mathrm{gap}(\widetilde{x}_T)$ после $T$~итераций выполнено
\begin{equation}\label{eq:viqa_rate}
  \mathrm{gap}(\widetilde{x}_T)
  = O\!\left(\frac{L_1 D^3}{T^{3/2}} + \frac{\delta D^2}{T}\right),
\end{equation}
где $D$~--- диаметр $\mathcal{X}$. Первое слагаемое~--- оптимальная
скорость точного метода второго порядка; второе~--- штраф за~неточность
якобиана, линейный по~$\delta$. В~частности, при $\delta\le L_1 D/\sqrt{T}$
штрафное слагаемое не~превосходит первого, и~метод сохраняет точную
скорость $O(L_1 D^3/T^{3/2})$. Тот~же вывод даёт и~более слабое
\emph{относительное} условие на~ошибку,
$\Norm{(\nabla F(v)-J(v))[x-v]}\le\tfrac{L_1}{2}\Norm{x-v}^2$, при котором
штрафное слагаемое поглощается кубическим членом~\eqref{eq:viqa_Omega}.
Тем самым любое количественное усиление оценки $\Norm{B_k-J}$~--- предмет
глав~\ref{ch:sp-broyden-theory} и~\ref{ch:symmetric}~--- напрямую
переносится в~оценку~\eqref{eq:viqa_rate} внешнего метода.


\section{Нижняя оценка и оптимальность}\label{sec:viqa_lower}

Совпадающая нижняя оценка (Теорема~4.3 из~\cite{viqa2024})
\begin{equation}\label{eq:viqa_lower}
  \mathrm{gap}(\widetilde{x}_T)
  = \Omega\!\left(\frac{\delta D^2}{T} + \frac{L_1 D^3}{T^{3/2}}\right)
\end{equation}
для класса методов дуально-экстраполяционного типа с~$\delta$-неточным
оракулом показывает, что зависимость~\eqref{eq:viqa_rate} от~$\delta$
неулучшаема: уровень неточности якобиана входит в~скорость по~существу,
а~не~как артефакт анализа. Следовательно, единственный рычаг ускорения
при~фиксированных $L_1, D, T$~--- уменьшение самого $\delta$, то~есть
качество квазиньютоновской аппроксимации якобиана.


\section{Рестарты: VIJI-Restarted при сильной монотонности}\label{sec:viqa_restart}

Если оператор дополнительно \emph{сильно монотонен},
$\langle F(x)-F(y),\,x-y\rangle\ge\mu\Norm{x-y}^2$ при $\mu>0$, то~базовую
схему ускоряет внешний рестарт-цикл VIJI-Restarted (Алгоритм~2
из~\cite{viqa2024}). Метод запускается $\lceil\log(D/\varepsilon)\rceil$
раз; на~$i$-й перезапуск VIJI выполняется
\begin{equation}\label{eq:viqa_restart_Ti}
  T_i = O\!\left(\max\Bigl\{
     \bigl(\tfrac{L_1 R_{i-1}}{\mu}\bigr)^{2/3},\ \tfrac{\delta}{\mu}
     \Bigr\}\right)
\end{equation}
итераций, после чего радиус локализации решения сокращается вдвое,
$R_i = R/2^{\,i}$. Суммарное число итераций до достижения
$\Norm{z_s-x^*}\le\varepsilon$ (Теорема~6.2 из~\cite{viqa2024}) составляет
\begin{equation}\label{eq:viqa_restart_rate}
  O\!\left(\Bigl(\tfrac{L_1 D}{\mu}\Bigr)^{2/3}
           + \Bigl(\tfrac{\delta}{\mu}+1\Bigr)
             \bigl\lceil\log(D/\varepsilon)\bigr\rceil\right).
\end{equation}
Здесь неточность якобиана входит \emph{аддитивно} и~лишь
с~логарифмическим множителем $\lceil\log(D/\varepsilon)\rceil$: первое
слагаемое~--- оптимальная второпорядковая сложность точного метода, второе
отражает $\delta$. Поэтому удержание $\delta$ малым~--- предмет
глав~\ref{ch:sp-broyden-theory} и~\ref{ch:symmetric}~--- напрямую
сокращает рестарт-сложность~\eqref{eq:viqa_restart_rate}, причём
из~\eqref{eq:viqa_restart_Ti} видно, что при $\delta\lesssim
(L_1 R_{i-1}/\mu)^{2/3}\mu$ штраф за~неточность вообще не~увеличивает длину
сегмента.


\section{Квазиньютоновская постановка}\label{sec:viqa_qn}

Оценка~\eqref{eq:viqa_rate} превращает задачу «удешевить метод второго
порядка» в~задачу «удержать $\delta$ малым дёшево», что и~есть
квазиньютоновская постановка. В~\cite{viqa2024} приближение $J(v)$
строится \emph{односекущими} обновлениями Бройдена с~ограниченной
памятью~--- L-Broyden,
\begin{equation}\label{eq:viqa_lbroyden}
  J_{i+1} = J_i + \frac{(y_i - J_i s_i)\,s_i^\top}{s_i^\top s_i},
  \qquad i = 0,\dots,m-1,
\end{equation}
и~его демпфированным вариантом; пары $(s_i,y_i)$ берутся либо из~истории
оператора ($s_i=z_{i+1}-z_i$, $y_i=F(z_{i+1})-F(z_i)$), либо из~якобиан-векторных
произведений $y_i=\nabla F(x)\,s_i$ на~случайных направлениях. Для них
доказана оценка $\delta\le(m{+}2)\,L_0$ (L-Broyden) и~$\delta\le 2L_0$
(демпфированный вариант)~\cite{viqa2024}. Здесь $L_0$~--- константа Липшица
\emph{оператора}~$F$, то~есть равномерная оценка нормы якобиана
$\Norm{\nabla F(\cdot)}_{\mathrm{op}}\le L_0$. Эти
оценки~--- априорный потолок: демпфирование удерживает
$\Norm{J}_{\mathrm{op}}\le L_0$, откуда
$\delta\le\Norm{\nabla F}_{\mathrm{op}}+\Norm{J}_{\mathrm{op}}\le 2L_0$
при~любом числе обновлений; качество же аппроксимации проявляется в~том,
что вблизи решения $\delta$ оказывается много меньше этого потолка. Существенно, что~\eqref{eq:viqa_lbroyden}~---
обновление ранга~1: на~каждом шаге сохраняется лишь текущее секущее
условие, то~есть тот самый ранг-1 проектор, расширение которого
до~ранга~$p{+}1$ составляет содержание главы~\ref{ch:sp-broyden-theory}.

Решение подзадачи~\eqref{eq:viqa_subcrit} с~низкоранговым $J$ дополнительно
удешевляется формулой Вудбери: при~ранге аппроксимации~$r$ стоимость
падает с~$O(n^3)$ до~$\widetilde{O}(r^2 n + r^3\log(D/\varepsilon))$. Это
прямо мотивирует малую память и~низкий ранг приближения~--- ту~же линию,
по~которой в~\S\ref{sec:highdim} строится limited-memory вариант~L-PB
с~окном ранга~$p{+}1$.

Обе линии работы относятся к~одной и~той же величине. Уровень неточности
$\delta$ из~\eqref{eq:viqa_delta} есть в~точности ошибка аппроксимации
якобиана: при~$J(v)=B_k$ выполнено
$\delta=\Norm{B_k-\nabla F(v)}_{\mathrm{op}}\le\Fro{B_k-\nabla F(v)}$~---
та~самая величина $\Norm{B_k-J}$, которую анализируют
главы~\ref{ch:sp-broyden-theory} и~\ref{ch:symmetric}. Различается \emph{вид}
доступного над~ней контроля. Односекущая формула~\eqref{eq:viqa_lbroyden}
даёт равномерный априорный потолок $\delta\le 2L_0$, не~зависящий ни~от
размерности, ни~от~траектории. Мультисекущее обновление удерживает
на~каждом шаге $p{+}1$ секущих условий вместо одного~--- строго больше
информации об~ошибке (следствие~\ref{cor:subspace-alignment})~--- и~даёт
для~$\Norm{B_k-J}$ оценку ограниченной деградации (теорема~\ref{thm:bd}; для
симметричного случая~--- композитное обновление SS-$\mathcal{U}$,
Лемма~\ref{lem:bd_composite}) с~константой
$C_{\mathrm{PB}} = L_1\,(1{+}C_s/2)\sqrt{p+1}\,\kappa_p$, зависящей
от~обусловленности окна~$\kappa_p$, а~\emph{не}~от размерности~$n$.

\section{Подстановка мультисекущей оценки}\label{sec:viqa_subst}

Глобальный потолок $\delta\le 2L_0$ односекущего L-Broyden не
убывает по~ходу итераций; напротив, мультисекущая оценка
теоремы~\ref{thm:bd}~--- \emph{локальная} рекуррентная оценка с~членом
деградации $C_{\mathrm{PB}}^2\rho_k^2$, действующая в~окрестности решения,
где $\rho_k=\max_{k-p\le j\le k}\Norm{x_j-x^*}$~--- радиус окна последних
шагов. В~режиме сходимости $\rho_k\to0$, и~штрафное слагаемое~$\delta$
становится убывающим. Это фиксирует следующее утверждение.

\begin{proposition}[Локальная подстановка мультисекущей оценки]\label{prop:viqa_subst}
Пусть в~методе~\S\ref{sec:viqa_method} якобиан аппроксимируется
мультисекущим обновлением Бройдена $J(v)=B_k$ (формула~\eqref{eq:sp_update}), а
траектория $\{x_k\}$ удержана в~окрестности решения, где применима
теорема~\ref{thm:bd}. Обозначим $E_k = B_k - \nabla F(x^*)$~--- ошибку
аппроксимации относительно якобиана в~решении (как в~теореме~\ref{thm:bd}).
Тогда уровень неточности~\eqref{eq:viqa_delta} в~текущей точке $v=x_k$
допускает оценку
\[
  \delta_k = \Norm{B_k-\nabla F(x_k)}_{\mathrm{op}}
           \le \Fro{E_k} + L_1\,\rho_k,
\]
по~неравенству $\Norm{\cdot}_{\mathrm{op}}\le\Fro{\cdot}$ и~$L_1$-гладкости
оператора, $\Norm{\nabla F(x^*)-\nabla F(x_k)}_{\mathrm{op}}\le L_1\Norm{x_k-x^*}
\le L_1\rho_k$ (\S\ref{sec:viqa_setup}). Существенный вклад мультисекущего
обновления~--- контроль ошибки \emph{на~подпространстве прошлых шагов}:
по~следствию~\ref{cor:subspace-alignment} $\Fro{E_k\,S_p} = O(\rho_k)$
с~константой $C_{\mathrm{PB}}=L_1(1{+}C_s/2)\sqrt{p+1}\,\kappa_p$, зависящей
от~обусловленности окна~$\kappa_p$, но~не~от размерности~$n$,~--- тогда как
односекущий L-Broyden удерживает лишь равномерный потолок $\delta\le 2L_0$,
от~$\rho_k$ не~убывающий. Подстановка этой подпространственной оценки
в~роль~$\delta$ отвечает локально убывающему по~$\rho_k$ штрафу
$O(\rho_k D^2/T)$ в~\eqref{eq:viqa_rate} вместо постоянного $O(L_0 D^2/T)$;
полная операторная норма~$\delta_k$ такого убывания, вообще говоря,
не~наследует~--- член $\Fro{E_k}$ по~теореме~\ref{thm:bd} остаётся
ограниченным, не~стремясь к~нулю.
\end{proposition}

\noindent
Утверждение~\ref{prop:viqa_subst} следует напрямую из~неравенства
$\Norm{\cdot}_{\mathrm{op}}\le\Fro{\cdot}$, $L_1$-гладкости оператора
и~следствия~\ref{cor:subspace-alignment};
оно сознательно \emph{локально} и~не~является замкнутой глобальной
скоростной теоремой: теорема~\ref{thm:bd}~--- рекуррентное неравенство,
действующее в~окрестности решения, тогда как
оценка~\eqref{eq:viqa_rate}~--- глобальная. Строгое согласование двух
уровней~--- вывод замкнутой скоростной оценки, в~которой $\delta_k$
мультисекущего обновления подставлено в~анализ
сходимости~\eqref{eq:viqa_rate} на~всей траектории,~--- остаётся
направлением дальнейшей работы.

\paragraph{Резюме.}\label{par:viqa-summary}
Глава помещает локальную теорию глав~\ref{ch:sp-broyden-theory}
и~\ref{ch:symmetric} в~контекст глобального метода второго порядка для
вариационных неравенств~\cite{viqa2024}. Развёрнута сама итерация дуальной
экстраполяции (\S\ref{sec:viqa_method}, Алгоритм~1 из~\cite{viqa2024})
и~её рестарт-вариант при~сильной монотонности
(\S\ref{sec:viqa_restart}, VIJI-Restarted): уровень неточности якобиана
$\delta=\Norm{\nabla F - J}$ входит в~скорость~\eqref{eq:viqa_rate}
линейно и~неулучшаемо (\S\ref{sec:viqa_lower}), а~в~рестарт-сложность
~\eqref{eq:viqa_restart_rate}~--- аддитивно-логарифмически. Удержание
$\delta$ возлагается на~квазиньютоновскую аппроксимацию: в~\cite{viqa2024}
эту роль играет односекущий L-Broyden с~потолком $\delta\le(m{+}2)L_0$;
мультисекущее обновление главы~\ref{ch:sp-broyden-theory} и~SS-$\mathcal{U}$
главы~\ref{ch:symmetric} дают для~той же величины на~исследованном
подпространстве локально убывающую оценку $O(\rho_k)$ с~геометрической
константой $\kappa_p$ вместо размерности~$n$
(утверждение~\ref{prop:viqa_subst}). Замкнутый глобальный вывод скоростной
оценки с~подставленным $\delta_k$ остаётся естественным направлением
дальнейшей работы.